\documentclass[11pt,letterpaper]{article}

\usepackage[margin=1in]{geometry}
\usepackage{amsmath}
\usepackage{graphicx}
\usepackage{hyperref}
\usepackage{booktabs}
\usepackage{array}
\usepackage{xcolor}
\usepackage{enumitem}
\usepackage{caption}
\usepackage{subcaption}
\usepackage{microtype}
\usepackage{parskip}
\usepackage{tikz}
\usetikzlibrary{arrows.meta, positioning, fit}

\hypersetup{
  colorlinks=true,
  linkcolor=blue!60!black,
  urlcolor=blue!60!black,
}

\title{\textbf{Brain Gene Expression Explorer}\\
  \large Interactive 3D Visualization of Spatially-Resolved\\
  Gene Expression Across a Whole-Brain Parcellation}

\author{Thiago Viegas\\
  \small New York University}

\date{\today}

% ─────────────────────────────────────────────────────────────────────────────
\begin{document}

\maketitle

\begin{abstract}
Spatially resolved gene expression is important for understanding how molecular
patterns vary across the human brain, but the outputs of transcriptomic alignment
models are difficult to evaluate when stored as dense numerical arrays across
many genes, subjects, and brain parcels.
We present \emph{Brain Gene Expression Explorer}, an interactive 3D visualization
tool developed to support ongoing model-development work in the NYU
Neuroinformatics Lab.
The tool visualizes gene expression predictions generated by three lab-developed
model variants: naive, DLAM, and PLAM.
Its main goal is to help developers visually compare each model's predictions
against available ground-truth measurements and inspect where the models agree,
disagree, or produce spatially structured errors. The application maps expression values onto a 156-parcel whole-brain atlas and
renders them in a browser-native WebGL viewer.
A four-panel layout simultaneously displays raw GTEx measurements, held-out
cross-validated reconstructions, dense whole-brain model predictions, and an
AHBA-based reference view.
This design allows developers to compare sparse measured data, model predictions,
and residual patterns within the same anatomical frame of reference.
A residual mode switches the reconstruction and whole-brain prediction panels
to signed error maps, making over-prediction and under-prediction visually
salient.
Developers can explore 88 high-variance genes across five donors and three model
variants with shared camera controls, a half-brain sagittal clip, parcel-level
tooltips, and independent opacity sliders for cortex, subcortex, and cerebellum.
A toggleable Model Analytics panel (V5) supplements the colour-based brain views
with position- and length-encoded statistical charts: a per-gene mean-expression
scatter and a ranked Pearson-$r$ bar chart for all genes provide a population-level
overview, while a colour-coded held-out parcel scatter and a paired horizontal bar
chart give gene-specific quantitative detail, offering an encoding channel that does
not depend on colour perception alone.
\end{abstract}

% ─────────────────────────────────────────────────────────────────────────────
\section{Introduction \& Motivation}
\label{sec:intro}

\subsection{Problem Definition}

The NYU Neuroinformatics Lab is developing machine-learning models that align
GTEx brain RNA-seq data with the Allen Human Brain Atlas (AHBA) to predict
spatial gene expression across brain parcels.
However, the models output dense numerical arrays across many genes, donors,
and regions, making it difficult to understand their behavior from tables or
summary metrics alone.

This project visualizes model predictions, ground-truth measurements, and
residual errors directly in anatomical space.
The goal is to help model developers compare the naive, DLAM, and PLAM models,
identify spatial patterns of agreement, failure, and bias, and evaluate how each
model extrapolates sparse measurements to the whole brain.
Beyond model debugging, the tool also helps users explore gene-expression
patterns themselves by showing how individual genes vary across anatomical
regions in a shared 3D brain view.
\subsection{Why It Matters}

Gene expression shapes brain function, development, and disease.
Spatially resolved expression maps can help researchers nominate candidate
genes for region-specific phenotypes, study transcriptomic gradients, and
generate hypotheses about how molecular patterns vary across the brain.
However, these patterns are difficult to interpret when they are stored only as
large arrays of values across genes, donors, and parcels.

This problem is also important from a machine-learning perspective.
In this project, the models are not only producing numerical estimates; they are
producing spatial hypotheses about where genes are expressed across the brain.
A model may achieve reasonable aggregate performance while still making
anatomically unrealistic predictions, smoothing away meaningful regional
differences, or performing poorly in specific structures such as the cerebellum
or subcortex.
Model developers therefore need a way to inspect predictions visually before
using them for downstream scientific analysis.

Existing neuroimaging viewers can handle continuous volumetric or surface data,
but they are not designed for rapidly switching between hundreds or thousands
of gene channels or comparing multiple model-variant outputs side by side.
This motivates an interactive visualization tool that connects gene-level
exploration with model debugging, allowing developers to compare predictions,
ground-truth measurements, and residual errors in a shared anatomical view.

\subsection{Target Users / Use Case}

The primary users are neuroinformatics researchers and graduate students working with the GTEx--AHBA integration pipeline at NYU.
Typical use cases include:
\begin{enumerate}[nosep]
  \item inspecting whether a model's predictions agree with sparse measured values for a specific gene and donor;
  \item comparing two model architectures, such as DLAM and PLAM, at a glance;
  \item identifying spatial residual patterns that may reveal model bias or
        anatomically meaningful differences;
  \item generating figures for presentations by exporting specific views.
\end{enumerate}

% ─────────────────────────────────────────────────────────────────────────────
\section{Data Description}
\label{sec:data}

\subsection{Datasets}

\paragraph{GTEx brain RNA-seq.}
Bulk RNA-seq count matrices for five post-mortem brain donors
(GTEX-1117F, GTEX-13OW8, GTEX-11DZ1, GTEX-1B996, GTEX-1JMPZ)
downloaded from the GTEx portal.
Each donor contributes between 8 and 12 anatomically distinct tissue samples,
each mapped to a parcel in the 4S156Parcels atlas.

\paragraph{Allen Human Brain Atlas (AHBA).}
Microarray expression data for approximately 900 tissue samples from 6 donors.
The upstream pipeline uses \texttt{abagen} to process and normalize these
data before spatial imputation.

\paragraph{4S156Parcels atlas.}
A whole-brain parcellation with 100 cortical and 56 subcortical regions
defined in MNI152NLin6Asym space.
Atlas metadata, including parcel labels, MNI centroids, and structure
assignments, are stored in a reformatted CSV.
The 3D NIfTI volume is used to extract subcortical meshes, while the
fsLR32k cortical surface files are used for cortical colouring.

\subsection{Preprocessing}

The upstream integration pipeline produces per-subject \texttt{.npz} archives
(one per model variant) that contain the following arrays, all aligned to
the canonical 150-parcel row order defined in \texttt{parcel\_order.json}
(alphabetical AHBA labels):

\begin{center}
\begin{tabular}{llll}
\toprule
Array key & Shape & Panel & Model-dep.\\
\midrule
\texttt{loro\_truth\_subject\_raw} & (150, $G$) & V1 & No \\
\texttt{loro\_eval\_mask}         & (150,)     & V1 mask & No \\
\texttt{loro\_fused\_subject\_h}  & (150, $G$) & V2 & Yes \\
\texttt{gtex\_mask}               & (150,)     & V2 mask & No \\
\texttt{fullfit\_subject\_h}      & (150, $G$) & V3 & Yes \\
\texttt{imputed\_mask}            & (150,)     & V3 mask & No \\
\texttt{gene\_names}              & (13618,)   & all & No \\
\bottomrule
\end{tabular}
\end{center}

\noindent $G = 13{,}618$ genes.
The viewer further filters to 88 high-variance genes listed in
\texttt{ahba\_100hvg.txt} to keep the dropdown usable.

\subsection{Challenges}

\paragraph{Atlas--data alignment.}
The 4S156Parcels atlas has 156 parcels; 6 lack AHBA microarray coverage
(\texttt{Cerebellar\_Region9}, \texttt{LH-MN}, \texttt{RH-EXA},
\texttt{RH-STH}, \texttt{RH-VeP}, \texttt{RH\_SomMot\_2}) and are
excluded from the \texttt{.npz}.
The application must align the 150-row data arrays to the 156-entry
atlas metadata without relying on integer index order.

\paragraph{Hemisphere convention.}
The upstream pipeline forces all sample MNI $x$-coordinates to be negative
(left hemisphere) before fitting.
As a result, some GTEx parcels are labelled as right-hemisphere structures
in the atlas even though the model treated them as left.
V1 and V2 apply a nearest-neighbour RH$\to$LH remap at render time;
V3 mirrors LH predictions onto their RH counterparts to produce a symmetric
bilateral view.

\paragraph{Mesh scale.}
Subcortical structures range from large nuclei, such as the putamen and
thalamus, to tiny parcels that cannot form a closed mesh at 1\,mm isotropic
resolution.
\texttt{LH-MN} and \texttt{RH-HN} were dropped from the mesh cache for this
reason; neither carries GTEx data.

% ─────────────────────────────────────────────────────────────────────────────
\section{Design \& Methodology}
\label{sec:design}

\subsection{Visualization Choices}

\paragraph{Four-panel side-by-side layout.}
The central design decision is to show all four data views simultaneously
on a shared camera.
This uses a \emph{juxtaposition} multiple-views strategy:
V1 shows sparse ground-truth measurements, V2 shows out-of-sample predictions
at those same locations, V3 shows dense whole-brain coverage, and V4 shows
the naive-model AHBA reference used as a fixed baseline.
Placing them side-by-side makes model accuracy and residual patterns visible
at a glance. A user can immediately compare V1 and V2, or V3 and V4,
without switching between separate screens.

\paragraph{3D surface encoding.}
We use a \emph{spatial position} primary encoding, with parcels rendered in
their anatomically correct MNI locations, plus a \emph{colour} secondary
encoding for expression level.
Position is the most direct visual channel for spatial anatomy; colour then
becomes the secondary channel for encoding continuous expression values once
position is already being used to represent anatomical location.

\paragraph{Colour.}
Two colormaps are used depending on context:
\begin{itemize}[nosep]
  \item \texttt{viridis} — sequential, for absolute expression values
        in V1, V2 Reconstruction, V3 Whole-brain Fit, and V4 AHBA Reference.
        It is perceptually uniform, greyscale-safe, and safe for the most
        common forms of colour-vision deficiency (deuteranopia and protanopia).
  \item \texttt{PuOr} (purple--orange) — diverging, for residual maps
        in V2 Error Analysis and V3 Residual.
        It is centred at zero so that over-prediction (orange) and
        under-prediction (purple) are immediately distinguishable.
        \texttt{PuOr} was chosen over red--blue alternatives because orange
        and purple remain perceptually distinct under both deuteranopia and
        protanopia, satisfying accessibility requirements.
\end{itemize}
The colormap range is auto-scaled \emph{per panel} by default; a Shared
mode locks all panels to a joint range, and a WBF mode locks all panels
to V3's range for direct cross-panel comparison.
Unmapped parcels, including parcels outside the panel's mask or parcels with
NaN values, are rendered in medium grey ($r{=}0.55, g{=}0.55, b{=}0.55$).
This darker neutral was chosen specifically to contrast with the near-white
centre of the \texttt{PuOr} diverging ramp, ensuring that parcels with no
data are not confused with near-zero residual values.

\paragraph{Model Analytics panel (V5).}
The three-dimensional brain panels encode expression through colour mapped onto
spatial position.
Colour is a powerful pre-attentive channel, but it has limitations as the
\emph{sole} encoding: it conveys one gene at a time, it does not support
quantitative comparison across many genes simultaneously, and readers with
colour-vision deficiency may find fine-grained distinctions difficult even with
accessible palettes.
To address this, we added a supplementary analytics panel, V5, that re-encodes
the same model outputs through \emph{position} (scatter plots) and \emph{length}
(bar charts) — visual channels that carry quantitative information more
precisely than colour and do not rely on colour discrimination.

V5 is organised in two rows, separated by a labelled divider, to distinguish
the scope of each view:

\begin{itemize}[nosep]
  \item \textbf{Row 1 — current gene.}
        Left: a colour-coded scatter plot of harmonized ground-truth versus
        LORO prediction at the held-out parcels, with each parcel labelled by
        region name.
        The identity line is drawn for reference; deviation from it represents
        prediction error.
        Right: a paired horizontal bar chart comparing harmonized truth and LORO
        prediction values parcel by parcel, sorted by truth magnitude, making
        systematic over- or under-prediction visible as consistent length
        discrepancies.

  \item \textbf{Row 2 — all genes.}
        Left: a global scatter in which each point represents one gene, plotted
        at its mean harmonized truth versus mean LORO prediction across held-out
        parcels for the active subject and model.
        The currently selected gene is highlighted in orange with a label,
        placing it in the context of the full gene population.
        Right: a ranked vertical bar chart showing Pearson $r$ for every gene in
        the 88-gene panel, sorted from lowest to highest, with the active gene
        highlighted.
        This chart makes it immediately visible which genes the model predicts
        well and which it predicts poorly, enabling the developer to judge
        whether the current gene is representative or an outlier.
\end{itemize}

Per-gene statistics (Pearson $r$, RMSE, mean truth, mean prediction) are
precomputed once at startup and recomputed only when the subject or model
changes, so switching genes regenerates only the chart image with no
additional array traversal.
V5 is hidden by default to preserve screen space and can be toggled from
the Analytics section of the side drawer.

\paragraph{Interaction design.}
Seven interaction types are provided:
\begin{enumerate}[nosep]
  \item \emph{Navigation} — mouse drag to orbit and scroll to zoom.
        All four panels move simultaneously through a shared VTK camera object.
  \item \emph{Selection} — gene, subject, and model dropdowns change the
        scalar array sliced from the in-memory NPZ.
  \item \emph{Filtering} — the half-brain toggle clips the scene at $x=0$
        and repositions the camera to a lateral view, exposing the
        mesial surface and deep structures.
  \item \emph{Error Analysis mode} — the Error Analysis toggle switches V2
        from absolute reconstruction to a prediction-error map, computed as
        LORO prediction minus harmonized ground truth. It simultaneously
        switches V3 to a whole-brain residual map, computed as WBF minus
        naive AHBA reference, for DLAM and PLAM models. The naive model
        stays in absolute mode because it is the reference itself.
  \item \emph{Click tooltip} — clicking any parcel displays a floating
        tooltip at the bottom centre of the viewport. In residual mode, the
        tooltip expands to three rows: Truth / Predicted / $\Delta$ for V2,
        and AHBA Reference / WBF / $\Delta$ for V3.
  \item \emph{Opacity adjustment} — three independent sliders control
        cortex (0--1), subcortex (0--1), and cerebellum (0--1) opacity,
        letting the user reveal or hide each compartment without discarding
        the others.
  \item \emph{Colour scale} — a three-way toggle, Local / Shared / WBF,
        controls whether panels auto-scale independently, share a joint
        range, or are all locked to V3's range.
\end{enumerate}

\subsection{Justification Using Course Concepts}

\paragraph{Perception \& pre-attentive features.}
Colour hue can be processed pre-attentively, making it useful for quickly
detecting regional expression differences across brain parcels.
Using a sequential rather than a categorical palette exploits the ordered
colour channel to encode an ordered quantity without requiring the viewer to
inspect every parcel individually.
Region boundaries are implied by the parcel mesh edges; no explicit contour is
needed because the visual system naturally groups nearby and continuous surfaces
as belonging together.

\paragraph{Visualizing machine-learning behavior.}
The dashboard design is motivated by the course idea that visualization can be
used not only to present final results, but also to inspect how a model is
behaving during development.
In this project, the models produce spatial predictions, so their behavior is
best evaluated spatially rather than only through aggregate metrics.
By placing ground-truth measurements, model reconstructions, whole-brain
predictions, and residual maps side by side, the visualization helps developers
ask whether the model's predictions make anatomical sense, where the model is
failing, and how each model extrapolates sparse measurements to the rest of the
brain.
The V5 analytics panel extends this to the gene population level: instead of
inspecting one gene at a time through colour, the per-gene $r$ bar chart and
global mean-expression scatter let developers immediately see whether the
current gene is among the best- or worst-predicted, and whether the model is
systematically biased across the expression range.
This bridges the gap between single-gene anatomical inspection and
aggregate model-performance evaluation.

\paragraph{Colour perception and accessibility.}
The medium-grey NaN colour ($r{=}0.55$) is dark enough to contrast clearly
with the near-white centre of the \texttt{PuOr} diverging ramp, preventing
near-zero residual values from being confused with unmapped parcels.
It is also distinct from the darker end of the \texttt{viridis} sequential
ramp, avoiding the salience bias that would arise if unmapped parcels were
rendered in a high-contrast colour such as black or red.
Both \texttt{viridis} and \texttt{PuOr} are recommended colourblind-safe
palettes: \texttt{viridis} is safe for deuteranopia and protanopia, and
\texttt{PuOr}'s orange--purple axis remains distinguishable under both
conditions. In addition, parcel-level click tooltips provide exact numerical
values, giving a non-colour secondary channel for residual magnitudes.

\subsection{Alternative Designs Considered}

\paragraph{2D flatmap.}
A cortical flatmap projection would eliminate occlusion and allow all parcels
to be visible simultaneously.
We rejected this because subcortical structures cannot be projected cleanly
onto a cortical surface flatmap, and the target users are more familiar with
3D anatomical views.

\paragraph{Heatmap matrix.}
A parcels $\times$ genes heatmap would support comparison of many genes at
once but would require the user to maintain a mental map from parcel label
to anatomical location.
The 3D brain surface leverages existing spatial knowledge and reduces the
cognitive load of label-to-anatomy lookup.

\paragraph{Plotly Dash.}
An earlier prototype used Plotly's \texttt{isosurface} trace.
It was discarded because loading a new gene required re-serialising the full
mesh to JSON on every update, making gene switching slow. Parcel-level colour
mapping was also fragile with the isosurface primitive.

% ─────────────────────────────────────────────────────────────────────────────
\section{System Description / Implementation}
\label{sec:system}

\subsection{Tools and Libraries}

\begin{center}
\begin{tabular}{lll}
\toprule
Library & Version & Role\\
\midrule
Python      & 3.10+  & application language \\
PyVista     & latest & VTK Python wrapper; mesh loading, scalar mapping \\
trame       & latest & Python-to-WebGL bridge; event loop and state management \\
trame-vtk   & latest & VTK render serialisation for trame \\
trame-vuetify & latest & Vuetify 3 / Vue 3 UI components defined in Python \\
wslink      & latest & WebSocket server backing trame \\
yabplot     & custom & fsLR32k surface loading and LUT parsing utilities \\
NumPy       & latest & array manipulation \\
matplotlib  & latest & colormap lookup table generation \\
\bottomrule
\end{tabular}
\end{center}

\subsection{Architecture Overview}

The application follows a \emph{server-side rendering} architecture:
all VTK geometry and scalar computations run on the Python server;
the browser client receives compressed differential frame updates
via WebSocket and displays them in a WebGL canvas.
This avoids the need to serialise large mesh data to JSON on every
interaction and keeps startup time low.

\begin{figure}[h]
  \centering
  \includegraphics[width=0.85\textwidth]{figures/architecture_figure.png}
  \caption{Application architecture. Input data are loaded on the Python server,
  where Trame manages application state and events, and PyVista/VTK performs
  geometry processing, scalar m
  apping, and rendering. The browser client sends
  interaction events through WebSocket and displays the returned rendered updates
  in a WebGL canvas.}
  \label{fig:arch}
\end{figure}

\subsection{Key Features}

\paragraph{Efficient gene switching.}
All model variants for the active subject are loaded into memory at
startup (\texttt{\_all\_npz}).
Changing the gene dropdown calls \texttt{np.where(gene\_names == gene)}
to retrieve the column index and slices the pre-loaded arrays in-place.
No disk I/O occurs during interaction, making gene switching instantaneous.

\paragraph{Shared camera.}
A single VTK \texttt{vtkCamera} object is shared across all four
plotter viewports.
Any mouse interaction, including orbit, zoom, or pan, applies to all panels
simultaneously, preserving spatial registration between views.
Programmatic camera changes, such as the half-brain lateral snap, are
applied via the VTK camera API directly, bypassing trame's state
serialisation path which can restore a stale camera position on
\texttt{view\_update()}.

\paragraph{Atlas alignment.}
\texttt{data\_loader.load\_alignment()} reads the atlas CSV and
\texttt{parcel\_order.json} and returns a DataFrame whose row order
matches the \texttt{.npz} arrays.
All downstream look-ups index by parcel label string, not by integer
position, making the alignment robust to reordering in either file.

\paragraph{RH$\to$LH remap.}
\texttt{build\_rh\_to\_lh\_map()} pre-computes a 69-entry
nearest-neighbour map from right-hemisphere parcel centroids to
their closest same-structure left-hemisphere counterparts.
V1 and V2 apply this map so that all data appear on the left
hemisphere, consistent with the upstream pipeline's spatial convention.

% ─────────────────────────────────────────────────────────────────────────────
\section{Results \& Insights}
\label{sec:results}

\subsection{Model Comparison}

The four-panel layout revealed systematic differences between
model variants and between predicted and measured values.

\paragraph{LORO cross-validation (V1 vs.\ V2).}
For most genes and donors, V2, which shows LORO reconstructions,
qualitatively matches V1, which shows raw GTEx measurements, at the
held-out parcels.
This confirms that the imputation models generalise to unseen locations.
Cerebellar parcels show the largest deviations. For example, for
\texttt{DPM1} in donor GTEX-13OW8, \texttt{Cerebellar\_Region1}
has V2 = 5.38, V3 = 5.09, and $\Delta = 0.30$.
This is consistent with the known difficulty of imputing expression in
the cerebellum from sparse donor data.

\paragraph{Model variant effects (V3).}
Switching the Model dropdown from \texttt{naive} to \texttt{dlam}
or \texttt{plam} in V3 reveals that the naive baseline produces
spatially inconsistent predictions in the right hemisphere, visible as
noise-like colour variation.
The DLAM and PLAM models produce smoother and more anatomically plausible
bilateral patterns.
This observation aligns with the design of the models: DLAM and PLAM
incorporate spatial regularisation terms that are absent from the naive
baseline.

\paragraph{Subcortical expression gradients.}
The cortex opacity slider, when lowered to approximately 0.3, exposes
subcortical structures that are otherwise partially hidden by the cortical
surface.
Several genes show strong subcortical--cortical contrasts.
For example, some genes appear more concentrated in striatal parcels than
in cortex, while others show stronger values in thalamic or cerebellar
regions.
These gradients are difficult to interpret from tables or heatmaps alone
and became more apparent in the 3D anatomical view.

\subsection{Residual Maps}

The Error Analysis mode, which switches V2 to a prediction-error map and V3
to a WBF-minus-reference residual map, proved useful for diagnosing
model-specific biases.
In V2, the prediction-error map
($\texttt{loro\_fused} - \texttt{loro\_truth}$)
highlights which held-out parcels a model over- or under-predicts.
Cerebellar parcels show some of the largest signed errors for the naive
baseline.

In V3, the WBF-minus-AHBA-Reference map isolates the spatial pattern of
improvement that DLAM and PLAM achieve over the naive model.
Positive regions indicate upward revision relative to the AHBA reference,
while negative regions indicate downward revision.
Because the naive model is the V4 reference, activating residual mode on
the naive subject correctly falls back to absolute display; there is no
residual to compute.

\subsection{Usability Observations}

The half-brain clip proved particularly useful for inspecting the
medial wall and cingulate cortex, which are occluded in the default
lateral view.
The three-way colour scale control, Local / Shared / WBF, addresses the
cross-panel comparison limitation noted in the previous prototype.
Shared mode locks all panels to a joint range so identical colours
represent identical expression values across panels, while Local mode
retains per-panel auto-scaling for maximum within-panel contrast.

\paragraph{Limited subject and gene-group support.}
The current version focuses on five GTEx donors and a curated set of 88
high-variance genes.
This is useful for early model debugging, but it limits broader exploration
across subjects and biological gene categories.
Future versions should support a larger set of donors and allow users to filter
or compare genes by predefined groups, pathways, or functional annotations.

% ─────────────────────────────────────────────────────────────────────────────
\section{Limitations \& Future Work}
\label{sec:limits}

\subsection{Current Limitations}

\paragraph{Cortical atlas rendering.}
The cortical surface currently renders as a solid grey shell because
the optional \texttt{wb\_command} pipeline step, which depends on
Connectome Workbench, was not completed in time for this submission.
Parcel-level cortical colour mapping requires projecting the volumetric
NIfTI atlas onto the fsLR32k surface; without this step, all cortical
expression is invisible.
Only the 54 subcortical VTK meshes carry parcel colours at this stage.

\paragraph{Per-panel colour scales in Local mode.}
Each panel auto-scales its colormap to the range of the currently
visible data when the Local scale mode is active.
This maximises within-panel contrast but means that identical colours
do not represent identical expression values across panels.
The Shared and WBF scale modes address this for V2/V3/V4 comparison,
but V1's sparse data, often only around 9 parcels, still benefits from
independent scaling.

\paragraph{No statistical overlay.}
The viewer shows point estimates only.
Uncertainty bounds, such as bootstrap confidence intervals on model
predictions, are available in the upstream pipeline output but are not
yet exposed in the UI.

\paragraph{Subject re-load latency.}
Switching subjects triggers a full NPZ reload from disk, usually around
1--2 seconds per model variant, with 3 variants loaded in parallel at startup.
For the current five-subject corpus this is acceptable, but it would become
a bottleneck with dozens of subjects.

\subsection{Future Work}

\begin{enumerate}[nosep]
  \item \textbf{Complete cortical atlas:} install Connectome Workbench
        and run \texttt{build\_atlas.py} to generate the fsLR32k parcel
        colour map, enabling full cortical visualisation.

  \item \textbf{Expand subject coverage:} add support for more GTEx donors
        and model outputs as they become available, allowing developers to
        compare whether model behavior is consistent across subjects.

  \item \textbf{Gene groups and functional categories:} allow users to explore
        predefined groups of genes, such as high-variance genes, pathway-related
        genes, or genes associated with specific biological functions. This would
        make the tool useful not only for single-gene inspection, but also for
        comparing broader transcriptomic patterns.

  \item \textbf{Additional coordinated views:} the V5 analytics panel already
        provides parcel-level scatter and bar charts linked to the active gene
        and subject.
        Future extensions could add gene-by-region heatmaps, subject comparison
        views, and bidirectional linking so that selecting a parcel in the 3D
        view highlights the corresponding point in V5, and vice versa.

  \item \textbf{Uncertainty visualisation:} render confidence intervals
        as a secondary scalar, such as hatching, secondary opacity, or
        stippling.

  \item \textbf{Gene search \& sorting:} replace the dropdown with a
        searchable list sorted by expression variance, differential expression
        across brain regions, or membership in a selected gene group.

  \item \textbf{Export:} add a screenshot button that captures the
        current view to a publication-quality PNG with a colour bar legend.

  \item \textbf{Multi-subject overlay:} show average expression across
        all available donors as an additional data view alongside individual
        subjects.
\end{enumerate}

% ─────────────────────────────────────────────────────────────────────────────
\section{Conclusion}
\label{sec:conclusion}

We presented the \emph{Brain Gene Expression Explorer}, an interactive
3D visualization tool that addresses the specific needs of researchers
working with spatially imputed gene expression data from the GTEx--AHBA
integration pipeline.
The key contributions are:

\begin{enumerate}[nosep]
  \item A four-panel side-by-side layout that makes model validation,
        cross-model comparison, and residual analysis visually immediate.
  \item A residual mode that switches V2 to a prediction-error map and V3
        to a whole-brain-fit-minus-AHBA-reference difference map, with
        per-parcel breakdown tooltips on click, supporting rapid diagnosis
        of model-specific spatial biases.
  \item An efficient architecture using server-side VTK, in-memory gene
        switching, and a shared camera to keep interactions responsive.
  \item A principled visual encoding using perceptually uniform sequential
        and diverging colourmaps, pre-attentive spatial mapping, and
        explicit handling of missing data.
  \item A Model Analytics panel (V5) that supplements colour-based brain
        views with position- and length-encoded statistical charts, providing
        a complementary encoding channel that does not depend on colour
        perception and enables population-level gene comparison alongside
        single-gene anatomical inspection.
  \item A hemisphere-remapping strategy that correctly handles the
        left-hemisphere convention of the upstream pipeline.
\end{enumerate}

The tool is actively used in the NYU Neuroinformatics Lab to diagnose
model behaviour and inspect expression patterns across brain regions,
demonstrating that targeted, domain-specific visualization tooling can
provide scientific value that general-purpose neuroimaging viewers
cannot easily replicate.

\end{document}